% -*- coding: utf-8 -*-

\begin{zhaiyao}
\begin{spacing}{1.5}
随着大语言模型从单轮问答逐步迈向跨会话、长时程交互，长期记忆系统已成为智能体保持连续性、个性化与任务一致性的关键基础设施。然而，现有评测大多停留在最终答题正确率层面，难以进一步回答错误究竟来源于记忆未被正确写入、运行时未能有效检索，还是生成阶段没有忠实利用证据。针对这一问题，本文面向长期记忆智能体评测场景，设计并实现了一套统一的细粒度归因评测框架。本文在调研大量记忆系统与评测文献的基础上，将异构记忆系统统一抽象为编码、检索、生成三个核心层面，并据此构建三个并行评测探针及一个最终归因代理，用于定位系统失效来源。围绕该框架，本文进一步设计统一数据契约、统一适配器协议、统一语义判分机制以及结构化结果落盘方式，使不同记忆系统能够在同一评测层下进行对齐分析。本文同时完成了框架工程实现，并预留了实验结果留白区域，以支持后续持续补充量化结果与跨系统对比。该工作为长期记忆系统的统一评测、问题定位与后续扩展提供了可复用基础。
\end{spacing}

\vspace{\baselineskip}
\noindent{\zihaoxiaosi\bfseries 关键词：}

\hspace*{2em}长时记忆智能体；细粒度评测；三层抽象；归因分析；适配器协议
\end{zhaiyao}

\begin{abstract}
\begin{spacing}{1.5}
As large language models move from single-turn question answering to cross-session and long-horizon interaction, long-term memory systems have become a critical infrastructure for maintaining continuity, personalization, and task consistency in intelligent agents. However, most existing evaluations remain at the level of final answer accuracy, and cannot further determine whether a failure is caused by incorrect memory encoding, ineffective runtime retrieval, or unfaithful answer generation. To address this problem, this thesis designs and implements a unified fine-grained attribution framework for evaluating long-term memory agents. Based on an extensive survey of memory systems and evaluation literature, heterogeneous memory systems are abstracted into three core layers, namely encoding, retrieval, and generation. On this basis, the framework introduces three parallel evaluation probes together with a final attribution agent for locating the source of failures. Around this design, the thesis further establishes a unified data contract, a unified adapter protocol, a unified semantic judgement mechanism, and a structured result persistence scheme, enabling different memory systems to be analyzed under the same evaluation layer. The engineering implementation of the framework is completed, while placeholders are reserved for pending experimental results and future cross-system comparisons. This work provides a reusable foundation for unified evaluation, failure diagnosis, and extensible benchmarking of long-term memory systems.
\end{spacing}

\vspace{\baselineskip}
\noindent{\zihaoxiaosi\bfseries Keyword}

\hspace*{2em}Long-term Memory Agents; Fine-grained Evaluation; Three-layer Abstraction; Attribution Analysis; Adapter Protocol
\end{abstract}
